\chapter{Conclusion and Future Work}

\section{Summary of Contributions}

This project presents the first empirically validated implementation of the DKLLMR23
laconic PSI protocol, which provides plausible post-quantum security from Ring-LWE.
The work spans two phases: a core cryptographic implementation phase establishing
correctness, and an engineering optimisation phase establishing scalability and
real-world viability.

The following concrete outcomes were achieved:

\begin{enumerate}
    \item \textbf{Core implementation (Phase I).} A full end-to-end implementation
          of Laconic Encryption --- comprising \textsf{KeyGen}, \textsf{LaconicEnc}, \textsf{Dec},
          and \textsf{WitGen} --- and the two-round PSI protocol, written in Go 1.24
          using the Lattigo library for Ring-LWE primitives. Correctness was validated
          across multiple test scenarios including basic intersection, empty intersection,
          full intersection, and repeated encryption with fresh randomness.

    \item \textbf{Memory reduction via batched witness generation (Phase II).} The
          central engineering contribution of this project. Batched witness generation
          reduces auxiliary space complexity from $S_{\text{aux}}^{\text{naive}} = \Theta(m \cdot D^2 \cdot \log q)$
          to $S_{\text{aux}}^{\text{batch}} = \Theta(B \cdot D^2 \cdot \log q) + \Theta(m \cdot D^2)$,
          with $B \ll m$. This theoretical result --- proved in Theorem~\ref{thm:batch-space} ---
          was validated empirically, reducing peak RAM from a fully-parallel worst
          case of $\approx 312$\,GB to a constant 6.5\,GB across all tested dataset
          sizes. The reduction factor of $\log_2 q \approx 58$ directly corresponds
          to the gadget decomposition expansion factor, confirming the theoretical analysis.

    \item \textbf{In-memory Merkle tree.} Loading the Merkle tree into RAM eliminates
          the SQLite I/O bottleneck that dominated Phase I witness generation,
          achieving a $21\times$ speedup in that phase. This engineering decision
          transforms witness fetching from an I/O-bound operation into a memory-bandwidth-bound
          one, which is the correct bottleneck to hit given the hardware characteristics
          of modern server systems.

    \item \textbf{Distributed execution on Google Compute Engine.} Horizontal
          partitioning of the server dataset across a 20--25 node GCE cluster reduces
          end-to-end execution time by approximately $15\times$, bringing large-scale
          screening runs within the bounds of overnight compliance batch windows. The
          distributed architecture achieves approximately 60\% parallel efficiency,
          with the remaining overhead attributable to coordination, network latency,
          and minor load imbalance --- all well-understood and quantified.

    \item \textbf{Scalability benchmark.} The first empirical performance characterisation
          of Ring-LWE-based laconic PSI, validated at $m = 10{,}000$ server records
          under 64-bit quantum security ($D = 256$), and correctness-validated at
          128-bit quantum security ($D = 2048$) on small datasets. These results
          establish the first measured data points for lattice-based laconic PSI,
          complementing the pairing-based empirical study of ALOS22~\cite{alos2022} and
          providing the research community with a concrete foundation for future
          post-quantum PSI development.

    \item \textbf{FLARE prototype.} A fully functional, GDPR-compliant sanctions
          screening web application demonstrating that the LE-PSI protocol can be
          deployed in a real-world financial compliance workflow. FLARE provides
          information-theoretic privacy guarantees that go beyond GDPR's contractual
          requirements: the sanctions authority is mathematically incapable of learning
          which customers the bank is screening, not merely contractually prohibited
          from doing so.
\end{enumerate}

% ─────────────────────────────────────────────────────────────────────────────
\section{Limitations}

This project makes several deliberate trade-offs in the interest of tractability
within the scope of an Honours project. These are stated explicitly and honestly:

\begin{enumerate}
    \item \textbf{Security level of large-scale benchmarks.} All large-scale
          performance results ($m = 10{,}000$) use $D = 256$, providing approximately
          64-bit quantum security --- equivalent to approximately 128-bit classical
          security, but below the NIST PQC standard of 128-bit quantum security~\cite{albrecht2015hardness}.
          This is a deliberate engineering trade-off: $D = 2048$ provides 128-bit quantum
          security but yields ciphertexts approximately $4\times$ larger and computations
          approximately $4\times$ slower. Scaling to 128-bit quantum security at
          $m = 10{,}000$ would require auxiliary space of approximately
          $B \times 280$\,MB (witness buffer) plus $\approx 100$\,GB (tree), which
          requires hardware beyond what was available, or a more constrained batch
          size. The path to 128-bit benchmarks is clearly identified in Section~7.3.

    \item \textbf{Semi-honest adversary model only.} The current implementation is
          secure only against semi-honest (honest-but-curious) adversaries --- parties
          who follow the protocol correctly but attempt to learn additional information
          from the transcript. It does not defend against malicious adversaries who
          may deviate from the protocol, submit malformed ciphertexts, or attempt
          active attacks. Extension to the malicious model requires zero-knowledge
          proofs of correct encryption from the client and correct intersection
          detection from the server, which is identified as future work.

    \item \textbf{Extrapolated 128-bit performance at $m = 10{,}000$.} While the $4\times$ computational overhead and $4\times$ memory expansion for 128-bit quantum security ($D = 2048$) were definitively established through direct empirical measurement on small datasets ($m \le 1{,}000$), executing this configuration at the full $m = 10{,}000$ scale on a single node remained intractable. Consequently, the performance projections for 128-bit security at $m = 10{,}000$ are mathematically extrapolated from the measured small-scale baseline overhead, rather than directly benchmarked end-to-end at boundary scale.

    \item \textbf{LWE Decryption Failure Probability at Small Datasets.} At very
          small dataset sizes ($m < 250$), the shallow Merkle tree limits the
          available entropy mixed into the ciphertext during gadget decomposition,
          occasionally causing LWE decryption failures due to insufficient error
          term ($\chi_e$) cancellation. This artefact results in reduced intersection
          accuracy (approximately 90\%) at $m = 100$. This is fundamentally a noise
          management issue, not a hash collision, and does not affect the
          results at the primary benchmark scale of $m \ge 1{,}000$ where tree depth
          provides sufficient distribution.

    \item \textbf{Communication extrapolation at large $m$.} Communication volume
          figures at $m = 10^6$ in Section~6.5 are extrapolated from the verified
          scaling model rather than directly measured. Direct measurement at
          $m = 10^6$ was not feasible within available computational resources.
\end{enumerate}

% ─────────────────────────────────────────────────────────────────────────────
\section{Future Work}

The results and limitations identified in this project suggest several high-impact
directions for future research and engineering:

\subsection{Scaling to 128-bit Quantum Security at Large Datasets}
The most immediate and direct next step is executing the $m = 10{,}000$ benchmark
at $D = 2048$ (128-bit quantum security). With a reduced batch size of $B = 5$
and a distributed setup of 3--5 GCE nodes, the per-node witness buffer drops to
approximately $5 \times 280$\,MB $\approx 1.4$\,GB, making the experiment feasible
on standard cloud hardware. This would establish the first 128-bit post-quantum
laconic PSI benchmark --- a significant milestone for the field.

\subsection{GPU Acceleration of NTT Operations}
NTT operations constitute approximately 25\% of execution time and are highly
data-parallel --- each polynomial coefficient is transformed independently. CUDA-based
implementations of Ring-LWE NTT (as demonstrated in cuFHE and HEAAN-GPU) have shown
$10\text{--}100\times$ speedups over CPU implementations. Porting the \texttt{pkg/matrix/}
NTT routines to a GPU kernel is therefore a high-impact optimisation that directly
targets the second-largest bottleneck identified in Section~6.6, without requiring
any changes to the cryptographic construction.

\subsection{Precise Security Parameter Estimation}
The current quantum security estimates use the BKZ cost model with standard assumptions.
Running the formal LWE Estimator~\cite{albrecht2015hardness} on the exact parameters
($n = 256, \sigma = 3.2, q \approx 2^{58}$) would yield a tighter, more defensible
security bound, and would also clarify the precise classical security equivalent ---
important for regulatory submissions where concrete security levels must be stated.

\subsection{Direct Comparison with ALOS22 Implementation}
The ALOS22 authors released their implementation at \url{https://eprint.iacr.org/2022/529}.
A head-to-head benchmark on identical hardware, with both implementations processing
the same dataset under matched conditions, would precisely quantify the performance
cost of replacing pairing arithmetic with Ring-LWE gadget decomposition. This would
provide the community with a definitive measurement of the post-quantum premium in
the laconic PSI setting.

\subsection{Malicious Security via Zero-Knowledge Proofs}
Extending the protocol to the malicious adversary model requires zero-knowledge proofs
of: (1) correct encryption from the client --- proving that each submitted ciphertext
was generated honestly using the server's public parameters; and (2) correct
intersection detection from the server --- proving that all decryption checks were
performed faithfully. Recent advances in efficient zkSNARKs for Ring-LWE operations
make this a tractable direction, though with significant additional computational overhead.

\subsection{Pipeline Parallelism in Distributed Execution}
The current distributed architecture transmits all client ciphertexts to all worker
nodes before intersection begins. A pipelined approach --- beginning intersection on
early partitions while later partitions are still being transmitted --- would reduce
the network latency component of the 40\% parallel efficiency loss identified in
Section~6.9, potentially pushing parallel efficiency above 75\%.

\subsection{Larger Dataset Validation}
Validating the distributed architecture at $m = 10^5$ and $m = 10^6$ records would
confirm the $O(n \log m)$ communication advantage of laconic PSI at scales where
the crossover with KKRT becomes practically significant (identified in Section~6.5 as
between $m = 10^4$ and $m = 10^5$). This would also validate the linear execution
time scaling observed at smaller scales and provide the first empirical data on
distributed laconic PSI at production-relevant sizes.

% ─────────────────────────────────────────────────────────────────────────────
\section{Conclusion}

This project demonstrates that the DKLLMR23 laconic PSI protocol --- previously a
purely theoretical construction --- can be made practical through careful memory
engineering and distributed system design. The central insight, proved in
Theorem~\ref{thm:batch-space} and validated empirically, is that auxiliary space complexity can be
reduced from $\Theta(m \cdot D^2 \cdot \log q)$ to $\Theta(m \cdot D^2)$ through
batched witness generation, transforming an impractical memory requirement into
one compatible with commodity hardware.

The work establishes several firsts: the first open implementation of DKLLMR23,
the first empirical performance characterisation of Ring-LWE-based laconic PSI,
the first distributed deployment of post-quantum laconic PSI, and the first
GDPR-compliance validation of laconic PSI in a financial screening context.

The performance overhead relative to classical alternatives --- approximately
$10{,}000\times$ slower than KKRT and $1{,}000\times$ slower than ALOS22 --- is
substantial and reflects the genuine cost of post-quantum security in the laconic
setting. This overhead is not an implementation deficiency; it is an inherent
consequence of Ring-LWE gadget decomposition and cannot be eliminated without
abandoning post-quantum security. What this project shows is that the overhead,
while large, is not prohibitive: with appropriate memory management and distributed
execution, the protocol is deployable on real hardware for real compliance workflows.

As quantum computing capabilities continue to advance and the timeline for cryptographically
relevant quantum computers becomes clearer, the case for post-quantum PSI in
long-lived financial data applications will only strengthen. This project provides
the empirical foundation and open implementation on which that transition can be built.
